\documentclass[11pt]{amsart}
\usepackage[T1]{fontenc}
\usepackage[utf8]{inputenc}
\usepackage{amssymb}
\usepackage[margin=1.05in]{geometry}
\usepackage[hidelinks]{hyperref}
\newcommand{\stamptime}{\the\numexpr\time/60\relax:\ifnum\numexpr\time-60*(\time/60)\relax<10 0\fi\the\numexpr\time-60*(\time/60)\relax}
\AddToHook{shipout/foreground}{\put(0.5\paperwidth,-\paperheight+0.55cm){\makebox[0pt]{\tiny\ttfamily unit/doc/\jobname.tex \ --- \ compiled \today, \stamptime}}}
\newtheorem{theorem}{Theorem}
\newtheorem{proposition}[theorem]{Proposition}
\newtheorem{lemma}[theorem]{Lemma}
\newtheorem{corollary}[theorem]{Corollary}
\theoremstyle{definition}
\newtheorem{definition}[theorem]{Definition}
\theoremstyle{remark}
\newtheorem{remark}[theorem]{Remark}
\newcommand{\Om}{\Omega}
\newcommand{\Map}{\mathrm{Map}_*}
\newcommand{\fib}{\mathrm{fib}}
\newcommand{\Tr}[1]{\mathsf{T}_{#1}}
\newcommand{\un}[1]{|#1|}
\newcommand{\code}[1]{\texttt{#1}}
\newcommand{\cT}{\mathsf{isTrunc}}
\newcommand{\cC}{\mathsf{isConn}}
\title{Finite truncations from replacement, and a 3-truncation of $S^2$}
\date{28 October 2026}
\begin{document}
\maketitle

\noindent\textbf{Purpose and status.} The informal proof for the
construction of finite truncations in Unit and for an inhabitant of the
interface \code{Trunc3} of \code{sag/pi3-s2-postnikov}
(see \code{doc/pi3-s2-postnikov.pdf}). It was written before the formal
proof, and each lemma below is a declaration
(\S\ref{sec:corr}). It was refereed independently before formalisation:
no mathematical error was found, and the corrections asked for (universe
levels, cited dependencies, order of generalisation before $J$) are
included below. The formal proof is checked: \code{SphereTruncation.SphereTrunc.trunc3}
is an inhabitant of \code{Trunc3}, from the sphere axiom and the
\code{Theory} assumptions, with no further assumption. Foundations: the
\code{Theory} assumptions (univalence, images, surjectivity of the image
maps; function extensionality derived); for \S\ref{sec:sphere} also the
sphere axiom. No other axiom. The construction is Rijke's (\emph{The join
construction}, \S7): $\|A\|_{n+1}$ is the image of
$a\mapsto\|a=-\|_n$. Here the join is not used. Replacement provides the
small image, and the inductive definition of connectedness replaces
the connectivity theorems of the paper.

\section{Conventions}

Unit indexing: $\cT\,m\,X$ is ``$X$ is $(m-2)$-truncated'' ($\cT\,0$ is
contractible, $\cT\,(m+1)\,X$ means all identity types are $\cT\,m$);
$\cC\,m\,X$ is ``$X$ is $(m-2)$-connected'' ($\cC\,0=\top$, and
$\cC\,(m+1)\,X$ means $\|X\|_{-1}$ and all identity types are $\cC\,m$).
Both are propositions (\code{trunc-prop}, \code{conn-prop}). A
proposition is $\cT\,m$ for every $m\ge1$ (\code{prop-trunc}), and
$\cC\,(m+1)$ implies $\cC\,m$ (\code{conn-down}).
$\fib_f(b)=\sum_a(f\,a=b)$. Everything lives in one universe $U=U_l$.

\section{Connected maps and truncated families}

\begin{lemma}[null types]\label{lem:null}
For every $m$, if $\cC\,m\,C$ and $\cT\,m\,D$, then
$\mathrm{const}:D\to(C\to D)$ is an equivalence.
\end{lemma}
\begin{proof}
Induction on $m$. For $m=0$, $D$ and $C\to D$ are contractible.
For $m+1$: being an equivalence is a proposition and $C$ is merely
inhabited, so let $c_0:C$. Put $\mathrm{ev}\,g=g\,c_0$; then
$\mathrm{ev}(\mathrm{const}\,d)=d$ by definition. For $g:C\to D$ and $c:C$,
the type $c=c_0$ is $\cC\,m$ and $g\,c=g\,c_0$ is $\cT\,m$. By the
induction hypothesis there is $h_c:g\,c=g\,c_0$ with
$\mathrm{const}\,h_c=\mathsf{ap}_g$, namely the inverse of $\mathrm{const}$
at $\mathsf{ap}_g:(c=c_0)\to(g\,c=g\,c_0)$. So
$\mathrm{const}(\mathrm{ev}\,g)=g$ by function extensionality with
$h_c^{-1}$. Hence $\mathrm{const}$ has the two-sided inverse
$\mathrm{ev}$.
\end{proof}

\begin{lemma}[elimination along a connected map]\label{lem:elim}
Let $f:A\to B$ with $\cC\,m\,(\fib_f\,b)$ for all $b$, and $P:B\to U$ with
$\cT\,m\,(P\,b)$ for all $b$. Then precomposition
$f^*:\prod_bP\,b\to\prod_aP(f\,a)$ is an equivalence.
\end{lemma}
\begin{proof}
Let $\Psi:\prod_b(\fib_f\,b\to P\,b)\to\prod_aP(f\,a)$, $\Psi\,t\,a=t\,(f\,a)\,(a,\mathsf{refl})$.
It has the inverse $\Phi\,s\,b\,(a,p)=p_*(s\,a)$: $\Psi\Phi=\mathrm{id}$ by
definition, $\Phi\Psi=\mathrm{id}$ by $J$ on $p$ (and function
extensionality). Let $K:\prod_bP\,b\to\prod_b(\fib_f\,b\to P\,b)$ be
$K\,s\,b=\mathrm{const}(s\,b)$. It is an equivalence: it is a $\Pi$ of
equivalences (Lemma~\ref{lem:null}), and a $\Pi$ of equivalences is an
equivalence by function extensionality. Finally $f^*=\Psi\circ K$ by
definition.
\end{proof}

In particular (constant $P$): for $\cT\,m\,Y$, $f^*:(B\to Y)\to(A\to Y)$ is
an equivalence; and maps $B\to Y$, or sections, that agree after $f^*$
are equal.

\begin{lemma}[fibres of $\mathsf{ap}$]\label{lem:fibap}
For $f:A\to B$, $x,y:A$ and $w:f\,x=f\,y$,
\[\fib_{\mathsf{ap}_f}(w)\ \simeq\ \big((x,w)=(y,\mathsf{refl})\big)\quad\text{in }\fib_f(f\,y).\]
Hence if all fibres of $f$ are $\cC\,(m+1)$, all fibres of
$\mathsf{ap}_f:(x=y)\to(f\,x=f\,y)$ are $\cC\,m$.
\end{lemma}
\begin{proof}
Paths in a $\Sigma$-type: $((x,w)=(y,\mathsf{refl}))\simeq\sum_{q:x=y}(q_*w=\mathsf{refl})$.
Moreover $q_*w=(\mathsf{ap}_fq)^{-1}\cdot w$ (\code{Groupoid.transport-fib}),
and $(\alpha^{-1}\cdot w=\mathsf{refl})\simeq(\alpha=w)$.
\end{proof}

\begin{lemma}[connectivity along a connected map]\label{lem:connup}
If all fibres of $f:A\to B$ are $\cC\,m$ and $\cC\,m\,A$, then $\cC\,m\,B$.
\end{lemma}
\begin{proof}
Induction on $m$; $m=0$ is trivial. For $m+1$: $B$ is merely inhabited
because $A$ is. For $b,b'$ the goal $\cC\,m(b=b')$ is a proposition and
the fibres over $b,b'$ are merely inhabited, so we may assume $b=f\,a$,
$b'=f\,a'$. Then apply the induction hypothesis to
$\mathsf{ap}_f:(a=a')\to(f\,a=f\,a')$: its fibres are $\cC\,m$ by
Lemma~\ref{lem:fibap}, and $a=a'$ is $\cC\,m$.
\end{proof}

\section{The truncations}

\begin{definition}[truncation structure at level $m$]
$\Tr m$ assigns to each $A:U$ a type $\Tr mA:U$ and a map
$\un-:A\to\Tr mA$ with $\cT\,m\,(\Tr mA)$ and $\cC\,m\,(\fib_{\un-}\,z)$
for all $z$.
\end{definition}

By Lemma~\ref{lem:elim} a truncation structure has dependent elimination
into $\cT\,m$ families, uniquely. For $m\ge1$ this includes families of
propositions.

\begin{lemma}[base, $m=1$]\label{lem:base}
The propositional truncation (\code{Truncation}) is a truncation
structure at level $1$: it is a proposition, and each fibre of its unit
is merely inhabited, by induction into propositions.
\end{lemma}

$\Tr0A=\top$ is also a truncation structure, but the step below is not
proved from $m=0$: it eliminates into families of propositions, which
are $\cT\,m$ only for $m\ge1$. So the recursion starts at $m=1$.

\paragraph{The step.} Let $\Tr m$ be a truncation structure, $m\ge1$, and $A:U$.
Let $f:A\to(A\to U)$, $f\,a\,c=\Tr m(a=c)$. The type $A\to U$ is
locally small, since $(g=h)\simeq\prod_c(g\,c\simeq h\,c)$ by function
extensionality and univalence (\code{Quotient.FunctionSmall}). So
replacement gives the image $I$ of $f$: a small type, a surjection
$q:A\to I$, an embedding $i:I\to(A\to U)$ and $t_a:i(q\,a)=f\,a$. Define
$\Tr{m+1}A=I$ and $\un a=q\,a$.

\begin{lemma}[Yoneda for truncated families]\label{lem:yoneda}
For $a:A$ and a family $Q:A\to U$ with $\cT\,m(Q\,c)$ for all $c$, the map
\[ \mathrm{Y}:\prod_c(\Tr m(a=c)\to Q\,c)\to Q\,a,\qquad
   \Phi\mapsto\Phi\,a\,\un{\mathsf{refl}} \]
is an equivalence. Its inverse sends $v$ to the family
$\Phi_v\,c=$ the extension along $\un-$ of $p\mapsto p_*v$.
\end{lemma}
\begin{proof}
$\mathrm{Y}$ is the composite of $\prod_c\un-^*$, which is an equivalence
for each $c$ by Lemma~\ref{lem:elim} with $Q\,c$ constant, and of the
evaluation at $\mathsf{refl}$, $\prod_c((a=c)\to Q\,c)\to Q\,a$. The
latter is an equivalence with inverse $v\mapsto(c,p)\mapsto p_*v$,
by $J$ on $p$ and function extensionality.
\end{proof}

\begin{lemma}[all such families are equivalences]\label{lem:alleq}
For $a,b:A$, every $\Phi:\prod_c(\Tr m(a=c)\to\Tr m(b=c))$ is a family of
equivalences.
\end{lemma}
\begin{proof}
By Lemma~\ref{lem:yoneda} with $Q\,c=\Tr m(b=c)$, $\Phi=\Phi_v$ with
$v=\mathrm{Y}\Phi:\Tr m(b=a)$ (the counit of the equivalence, a path,
along which the property is transported). So it suffices to prove, for
all $a$ and all $v:\Tr m(b=a)$, that $\Phi^{(a)}_v\,c$ is an equivalence
for all $c$. This is a proposition in $v$, so by elimination (here
$m\ge1$ is used) we may take $v=\un r$ with $r:b=a$. By $J$ on $r$, with
$a$ free, take $a\equiv b$ and $r=\mathsf{refl}$. Then $\Phi_{\un{\mathsf{refl}}}\,c$ extends
$p\mapsto p_*\un{\mathsf{refl}}=\un p$ (by $J$ on $p$), as the identity
of $\Tr m(b=c)$ does. By uniqueness of extensions it is the identity, an
equivalence.
\end{proof}

\begin{proposition}[identity types]\label{prop:ident}
$(q\,a=q\,b)\ \simeq\ \Tr m(b=a)$ for all $a,b:A$.
\end{proposition}
\begin{proof}
$(q\,a=q\,b)\simeq(i(q\,a)=i(q\,b))$ because $i$ is an embedding;
$\simeq(f\,a=f\,b)$ by conjugation with $t_a,t_b$; $\simeq\prod_c(f\,a\,c\simeq f\,b\,c)$
(\code{FunctionSmall}); $\simeq\prod_c(f\,a\,c\to f\,b\,c)$. The last step
holds because $\prod_c(X_c\simeq Y_c)$ is
$\sum_{\Phi}\prod_c\mathsf{isEquiv}(\Phi\,c)$, a $\Sigma$ of
propositions that all hold (Lemma~\ref{lem:alleq}). Finally
$\simeq\Tr m(b=a)$ by Lemma~\ref{lem:yoneda}.
\end{proof}

\begin{proposition}[level]\label{prop:level}
$\cT\,(m+1)\,I$.
\end{proposition}
\begin{proof}
For $z,w:I$ the goal $\cT\,m(z=w)$ is a proposition. Since $q$ is
surjective (induction into propositions) we may assume $z=q\,a$, $w=q\,b$.
Then use Proposition~\ref{prop:ident} and $\cT\,m(\Tr m(b=a))$.
\end{proof}

\begin{proposition}[connectivity of the unit]\label{prop:conn}
$\cC\,(m+1)\,(\fib_q\,z)$ for all $z:I$.
\end{proposition}
\begin{proof}
A proposition in $z$; so we may assume $z=q\,b$. By
Proposition~\ref{prop:ident} (a $\Sigma$ of equivalences),
$\fib_q(q\,b)\simeq S_b:=\sum_a\Tr m(b=a)$, and $\cC$ is invariant under
equivalence. $S_b$ is merely inhabited by $(b,\un{\mathsf{refl}})$. For
$(a,u),(a',u'):S_b$ the goal $\cC\,m((a,u)=(a',u'))$ is a proposition in
$u$. By elimination take $u=\un p$, $p:b=a$, and by $J$ on $p$, with
$a$ free, take $(a,u)=(b,\un{\mathsf{refl}})$. Then
\[ ((b,\un{\mathsf{refl}})=(a',u'))\ \simeq\ \textstyle\sum_{r:b=a'}(r_*\un{\mathsf{refl}}=u')
   \ \simeq\ \sum_{r:b=a'}(\un r=u')\ =\ \fib_{\un-}(u'), \]
by paths in a $\Sigma$-type and $r_*\un{\mathsf{refl}}=\un r$ ($J$ on
$r$). This is $\cC\,m$ because $\Tr m$ is a truncation structure.
\end{proof}

\begin{theorem}\label{thm:trunc}
For every $m\ge1$ there is a truncation structure $\Tr m$, by recursion on
$m$ from Lemma~\ref{lem:base} and the step
(Propositions~\ref{prop:level},~\ref{prop:conn}).
\end{theorem}

\section{Pointed universal property and $\pi_3$}

Let $\eta=(\un-,\mathsf{refl}):(X,x_0)\to_*(\Tr mX,\un{x_0})$.

\begin{proposition}[pointed universal property]\label{prop:ptup}
For $\cT\,m\,Y$, $\eta^*:\Map(\Tr mX,Y)\to\Map(X,Y)$ is an equivalence.
\end{proposition}
\begin{proof}
$\Map(\Tr mX,Y)=\sum_g(g\un{x_0}=y_0)$, and $(\un-^*g)\,x_0$ is
$g\un{x_0}$ by definition. So the $\Sigma$ over the equivalence $\un-^*$
of bases (Lemma~\ref{lem:elim}) is an equivalence onto
$\sum_h(h\,x_0=y_0)$. It agrees with $\eta^*$ up to
$\mathsf{refl}\cdot w=w$ in the pointing.
\end{proof}

\begin{lemma}[loopings of connected maps]\label{lem:loopconn}
If $F:(Y,y_0)\to_*(Z,z_0)$ has all fibres $\cC\,(n+1)$, then
$\Om F:\Om Y\to\Om Z$ has all fibres $\cC\,n$.
\end{lemma}
\begin{proof}
$\Om F$ is homotopic to $q\mapsto F_0^{-1}\cdot\mathsf{ap}_Fq\cdot F_0$
(\code{loop-ap}); by $J$ on the pointing $F_0$ we may take
$F_0=\mathsf{refl}$, and then $\Om F$ is $\mathsf{ap}_F$ up to
$\mathsf{refl}\cdot-$. Here $\Om F$ is the pointed map
\code{Maps.loop F}, and $\Om^3\eta$ below is \code{Maps.loop (loop2 eta)},
the map that \code{PiOne.Induced} uses for $\pi_3\eta$. Fibres of homotopic maps are equivalent
(\code{fib-htpy}). Conclude by Lemma~\ref{lem:fibap}.
\end{proof}

\begin{proposition}[$\pi_3$]\label{prop:pi3}
For $m=5$ (3-types), $\pi_3\eta:\pi_3X\to\pi_3(\Tr5X)$ is bijective.
\end{proposition}
\begin{proof}
The fibres of $\un-$ are $\cC\,5$. By Lemma~\ref{lem:loopconn} three
times, $\Om^3\eta$ has fibres $\cC\,2$: merely inhabited, with any two
points merely equal. $\pi_3\eta$ is the map induced by $\Om^3\eta$ on the
quotients of $\Om^3X$ and $\Om^3\Tr5X$ by mere equality. It is
surjective: by induction on the class and the mere inhabitation of
fibres. It is injective: if $\Om^3\eta\,q$ and $\Om^3\eta\,q'$ are merely
equal (effectiveness of the quotient by mere equality,
\code{Quotient.effective}; both classes by induction into a
proposition), then $(q,p)$ and $(q',\mathsf{refl})$ lie in one fibre, so they are
merely equal, so $q$ and $q'$ are. A map of sets with propositional,
merely inhabited fibres has contractible fibres, so it is an equivalence.
No choice is used.
\end{proof}

\section{The 3-truncation of $S^2$}\label{sec:sphere}

\begin{lemma}\label{lem:s2conn}
$\cC\,3\,S^2$ (simply connected).
\end{lemma}
\begin{proof}
Let $V=\Tr3S^2$, a 1-type. (a) For any 1-type $Y$, evaluation at
$\mathsf{base}$, $(S^2\to Y)\to Y$, is an equivalence: it is the first
projection after the axiom $\mathrm{ev}_2$, and the fibres
$\Om^2(Y,y)$ of the projection are contractible: $y=y$ is
$\cT\,2$, so $\Om^2(Y,y)$ is a proposition, inhabited by $\mathsf{refl}$. (b) The maps $\mathrm{id}$ and
$\mathrm{const}\,\un{\mathsf{base}}$ of $V$ agree after $\un-^*$. Indeed
$\un-$ and $\mathrm{const}\,\un{\mathsf{base}}$ agree at $\mathsf{base}$
by definition, so they are equal by (a) with $Y=V$. By the universal
property of $V$ (Lemma~\ref{lem:elim}), $\mathrm{id}=\mathrm{const}\,\un{\mathsf{base}}$,
so $V$ is contractible. (c) Then $\fib_{\un-}(\un{\mathsf{base}})\simeq S^2$:
a $\Sigma$ of contractible types $(\un s=\un{\mathsf{base}})$. The fibre
is $\cC\,3$, hence so is $S^2$.
\end{proof}

\begin{theorem}[an inhabitant of \code{Trunc3}]
Take $T=\Tr5S^2$, $t_0=\un{\mathsf{base}}$ and $\eta$ as above. Then
$\cT\,5\,T$ (Theorem~\ref{thm:trunc}); $\cC\,3\,T$ by
Lemma~\ref{lem:connup}, since the fibres of $\un-$ are $\cC\,5$, hence
$\cC\,3$ (\code{conn-down}), and $S^2$ is $\cC\,3$ (Lemma~\ref{lem:s2conn}); \code{up} is
Proposition~\ref{prop:ptup}; \code{pi3-eta} is Proposition~\ref{prop:pi3}.
\end{theorem}

\section{Universes}

All the constructed types are in $U_l$. The image is taken with
codomain $A\to U_l:U_{l+1}$, which is locally small: its identity types
are \emph{equivalent} to types of $U_l$ (\code{Small} is a type with an
equivalence). This uses function extensionality at $U_{l+1}$ (in
\code{Quotient.FunctionSmall}), derived from univalence at $l+1$, and
univalence at $l$. Replacement is used as \code{Images} and
\code{Surjectivity} at $(l,l+1)$, eliminating into $U_l$. These are
instances of the \code{Theory} parameters, which are level-polymorphic. $S^2$ and its
axiom are at level $l$, and the targets $Y$ of its universal property are
in $U_l$. In particular $V$ and $T$ are, as is the target of every
elimination.

\section{Correspondence with the formal declarations}\label{sec:corr}

\begin{tabular}{ll}
Lemma~\ref{lem:null} & \code{ConnectedMaps.null} \\
Lemma~\ref{lem:elim} & \code{ConnectedMaps.Elim}: \code{Psi}, \code{Phi}, \code{K}, \code{precomp-def}, \code{equiv}, \code{ext}, \code{ext-beta}, \code{unique}\\
Lemma~\ref{lem:fibap} & \code{ConnectedMaps.fib-ap}, \code{fib-ap-conn}\\
Lemma~\ref{lem:connup} & \code{ConnectedMaps.conn-up}\\
truncation structure & \code{Truncations.TruncOp}\\
Lemma~\ref{lem:base} & \code{Truncations.base}\\
Lemma~\ref{lem:yoneda} & \code{Truncations.Step.Yon}: \code{Y}, \code{equiv}, \code{Phi-of}, \code{Phi-of-Y}\\
Lemma~\ref{lem:alleq} & \code{Step.At.AllEq.all-equiv} (\code{id-at}, \code{good-unit}, \code{good})\\
Proposition~\ref{prop:ident} & \code{Step.At.Ident}: \code{e1}--\code{e4}, \code{ident}\\
Proposition~\ref{prop:level} & \code{Step.At.trunc'}\\
Proposition~\ref{prop:conn} & \code{Step.At.ConnS}, \code{Step.At.conn'}\\
Theorem~\ref{thm:trunc} & \code{Truncations.TO} (locked)\\
Proposition~\ref{prop:ptup} & \code{TruncationProps.PtUP.up}\\
Lemma~\ref{lem:loopconn} & \code{TruncationProps.loop-conn}\\
Proposition~\ref{prop:pi3} & \code{TruncationProps.ConnPi1.equiv}\\
Lemma~\ref{lem:s2conn} & \code{SphereTruncation.SphereTrunc}: \code{ev-base-equiv}, \code{V-contr}, \code{s2-conn}\\
the inhabitant of \code{Trunc3} & \code{SphereTrunc.trunc3}
\end{tabular}

\medskip
\noindent\emph{Checking against the text} (\code{seg/WRITING.md}).
\code{truncation.opens} declares, for each of the 94 declarations of
these four files, the definitions its step uses. These are the notions of
its statement and the components named in the text. The uniform checker
(\code{--opens}) reports nothing else opened, except in three places:
(i) \code{ConnPi1} opens the carriers of the $\pi_1$ quotient
(\code{PiOne.grp}, \code{QGrp.G}, \code{Quot.Q}, \code{image}), an interface
issue of the shared library, which was not edited; (ii) \code{c2} opens
\code{Maps.D}: the base point of a looping is $\mathsf{refl}$ by
computation; (iii) \code{V-contr}, \code{c5}, \code{connT} report
\code{TO} and the image of the step as opened, although \code{TO} is
locked. A test that needs \code{TO} to unfold is rejected by both
checkers, so no conversion relies on it. The cause of this report was
not determined. \code{TO}, \code{null}, \code{Elim.equiv} and \code{Elim.ext}
are locked after their characterising lemmas. Every declaration checks
in under 0.01\,s.

\end{document}
